\begin{abstract}
A procedure can be correct where it was learned and incomplete where it is reused. We study coding agents supplied with procedural memory in target tasks that change a security relevant source assumption. Thirteen synthetic families cross four memory conditions and two repetitions in six model and agent configurations, yielding \RecordedRuns{} recorded runs and \ValidRuns{} technically valid outcomes. The primary endpoint is functional completion that fails a focal security witness. Correct memory changes this endpoint by $\PrimaryURangeLow$ to $\PrimaryURangeHigh$ percentage points across configurations; the experiment does not establish a general harmful memory effect. Lower failure rates sometimes accompany lower functionality. A generic applicability reminder reduces the endpoint descriptively in the three Codex configurations, a direction that survives all assignments of unresolved outcomes in the recorded cohort. Two isolated LLM annotation passes over a sample of 52 runs document procedure retention, protection, explicit assumption recognition, and non-completion separately. Protections often appear without preceding explicit recognition in the visible trace. Together, the outcome decomposition and implementation evidence show how non-completion, witness scope, and distinct implementation mechanisms change the interpretation of apparent memory transfer successes and failures.
\end{abstract}
